% Isolated candidate; no canonical or shared build input is changed.
\subsection{Odd polarizations: Cartier surjectivity and generic evaluation}

Let \(A\) be an abelian surface over an algebraically closed field
\(k\) of characteristic two. Let \(L\) be ample and globally
generated, and suppose
\[
 N=\chi(L)=h^0(A,L)>1\quad\hbox{is odd}.
 \tag{DE-OE1}
\]
Put \(B=A^{(2)}\), let \(F:A\to B\) be relative Frobenius,
and write
\[
 E=F_*L,\qquad H=L^{(2)},\qquad W=H^0(B,E)=H^0(A,L).
 \tag{DE-OE2}
\]
Here \(H\) is the Frobenius twist of the line bundle, not its
square on the original surface: \(F^*H\simeq L^2\).
In particular \(H\) is ample and globally generated,
\(H^2=2N\), and numerically \(c_1(E)=2c_1(H)\).
The stability assertion of Lemma~\ref{lem:DE-odd-line-Frobenius-stable}
therefore gives \(\mu_H(E)=N\) and strict stability.

\begin{lemma}[A scaled divisor bound and the theta weights]
\label{lem:DE-odd-theta-section-degree}
\status{Candidate; complete independent audits pending.}
If a line bundle \(D\) on \(B\) has \(h^0(D)\ge N\), then
\[
                         c_1(D)\cdot H\ge2N.
 \tag{DE-OE3}
\]
Let \(\mathcal G(L)\) be the theta group. Every nonzero
finite-dimensional representation of central weight \(r\), with
\(\gcd(r,N)=1\), has dimension divisible by \(N\).
The induced action on \(E\) has central weight one and covers
translation by \(F(x)\) for \(x\in K(L)\); that on \(H\)
has central weight two. On the common finite translation subgroup,
\[
                 e_H(Fx,Fy)=e_L(x,y)^2.
 \tag{DE-OE4}
\]
\end{lemma}
\begin{proof}
An effective divisor on an abelian surface is nef, by translating
it to meet any specified curve properly. If \(D^2>0\), it is
ample: a curve orthogonal to it would have negative square by
Hodge index, whereas every curve on an abelian surface is nef.
Ample-line cohomology vanishing and Riemann--Roch give
\(D^2=2h^0(D)\ge2N\). Hodge index with \(H^2=2N\)
then gives (DE-OE3).

If \(D^2=0\), the complete effective-divisor argument in
Lemma~\ref{lem:DE-three-section-divisor-degree} applies without
the number three. Each component is a smooth elliptic translate,
zero pairwise intersections put them in one elliptic quotient
\(q:B\to C\), and \(D=q^*D_0\) with
\(\deg D_0=h^0(D)\ge N\). Global generation of \(H\)
implies degree at least two on every elliptic fibre. Hence
\(D\cdot H\ge2h^0(D)\ge2N\). The case \(D=0\)
cannot occur because \(N>1\).

The theta commutator is nondegenerate and \(K(L)\) has order
\(N^2\). Its odd order permits a Lagrangian decomposition
\(K(L)=K_1\times K_2\), where both factors have order \(N\),
and lifts to level subgroups. These are the full theta-group facts
of \citep[Proposition~8.19, Corollary~8.20, Lemma~8.24 and
Theorem~8.28 with proofs, pp.~119--123]{EdixhovenVanDerGeerMoonenAV}.
For completeness the needed dimension assertion is direct.
Decompose a representation under the first level subgroup into
character spaces. The other subgroup permutes those spaces by
multiplication with the characters \(e_L(-,y)^r\).
Because raising to \(r\) is invertible on these odd finite
groups, this is a simply transitive permutation of all \(N\)
characters. All their multiplicities are equal. A nonzero
representation therefore has dimension a positive multiple of
\(N\). In particular a representation of dimension \(N\)
and such a weight is irreducible.

A theta operator on \(L\) pushes forward to an operator on
\(E\) over the translation by \(F(x)\), with the same scalar
action. Its relative Frobenius twist acts on \(H\) with the
squared scalar action. Taking the commutator gives (DE-OE4).
Since \(K(L)\) has odd order, \(F\) identifies it with
\(K(H)\). All actions here are actions of actual theta groups;
no arbitrary projective lift is being chosen on a quotient line.
\end{proof}

\begin{lemma}[The evaluation saturation has rank at least three]
\label{lem:DE-odd-Cartier-evaluation-three}
\status{Candidate; complete independent audits pending.}
Let \(Q\subset E\) be the saturation of the evaluation image
\(W\otimes\mathcal O_B\to E\). Its rank is at least three.
\end{lemma}
\begin{proof}
It is positive because \(W\ne0\). The image and its saturation
are theta invariant. If its rank is \(q\in\{1,2\}\), wedge
evaluation gives a nonzero equivariant map
\[
          \bigwedge^q W\longrightarrow H^0(B,\det Q).
 \tag{DE-OE5}
\]
Its image has central weight \(q\). Since \(N\) is odd,
the preceding lemma implies \(h^0(\det Q)\ge N\).
It follows that \(c_1(Q)\cdot H\ge2N\).
On the other hand, strict stability of the rank-four bundle
\(E\) gives \(c_1(Q)\cdot H<qN\le2N\), a contradiction.
The determinant is the reflexive determinant; the evaluation
sections extend to it across the finite singular set, so (DE-OE5)
is an actual global map.
\end{proof}

\begin{theorem}[Surjectivity of odd-polarization Cartier multiplication]
\label{thm:DE-odd-polarization-Cartier-surjection}
\status{Candidate; complete independent audits pending.}
If \(A\) is nonordinary, relative Cartier multiplication induces
a surjection
\[
 \beta:\bigwedge^2 H^0(A,L)\longrightarrow H^0(B,H).
 \tag{DE-OE6}
\]
\end{theorem}
\begin{proof}
After choosing an invariant volume form and its Frobenius twist,
finite-flat duality and projection formula give the perfect
pairing \(E\otimes E\to H\) used in
Lemma~\ref{lem:DE-primitive-one-three-Cartier}.
It is alternating: on a local square, relative Cartier is
the pulled-out square factor times Cartier of the invariant
top form, and that latter scalar vanishes on a nonordinary
abelian surface. This establishes alternation on the bundle,
not only on global sections.

If \(\beta=0\), the generic evaluation subspace in the
four-dimensional symplectic fibre of \(E\) is isotropic,
so its dimension is at most two. This contradicts the preceding
lemma. Thus \(\beta\ne0\). The map is theta equivariant
and its target has central weight two and dimension \(N\).
The weight argument of Lemma~\ref{lem:DE-odd-theta-section-degree}
makes that target irreducible. Its nonzero invariant image is
therefore the whole target. This proves (DE-OE6).
\end{proof}

\begin{proposition}[Generic generation when the Euler characteristic is prime to six]
\label{prop:DE-prime-to-six-evaluation-four}
\status{Candidate; complete independent audits pending.}
If, in addition, \(3\nmid N\), the evaluation map
\(W\otimes\mathcal O_B\to F_*L\) is surjective at the
generic point. This assertion does not require nonordinarity.
\end{proposition}
\begin{proof}
The only remaining rank to exclude is three. Suppose the
evaluation saturation \(Q\) has rank three. Put
\(T=E/Q\), a torsion-free sheaf of rank one, and
\(J=T^{**}\), a line bundle. Saturation and reflexive hull
are functorial for the theta action, so \(J\) has an actual
\(\mathcal G(L)\)-action of central weight one, over translation
by \(F(x)\). Also \(D=\det Q\) has central weight three,
and wedge evaluation has nonzero image in \(H^0(D)\).
The coprimality assumption thus gives \(h^0(D)\ge N\).

We first prove that \(J\) is ample, rather than assuming
that a rank-one quotient of \(E\) is automatically ample.
The canonical filtration in
Lemma~\ref{lem:DE-odd-line-Frobenius-stable} makes
\(U=F^*E\otimes L^{-1}\) unipotent. Pulling back the quotient
and then passing to the double dual gives a map
\[
                   U\longrightarrow F^*J\otimes L^{-1}
 \tag{DE-OE7}
\]
which is surjective outside a finite set. For the normalization
\(\widetilde C\) of any integral curve on \(A\), its restriction
is generically surjective. The image is a line bundle of the form
\((F^*J\otimes L^{-1})|_{\widetilde C}(-Z_C)\),
where \(Z_C\) is effective. The restriction of \(U\) is an
extension of trivial bundles and is semistable of degree zero.
Every line-bundle quotient has nonnegative degree. Therefore
\(\deg((F^*J\otimes L^{-1})|_{\widetilde C})\ge0\).
The line \(F^*J\otimes L^{-1}\) is nef on \(A\).
Tensoring it with the ample line \(L\) shows that \(F^*J\)
is ample; finite surjective pullback detects ampleness, so \(J\)
is ample as claimed.

Hence \(h^0(J)=\chi(J)>0\). Its central weight is one, and
the theta dimension lemma gives \(h^0(J)\ge N\).
Apply (DE-OE3) to both \(D\) and \(J\). The finite-length
quotient \(J/T\) changes no first Chern class, so
\[
  D\cdot H\ge2N,\qquad J\cdot H\ge2N,\qquad
  (D+J)\cdot H=c_1(E)\cdot H=4N.
 \tag{DE-OE8}
\]
Both intersections equal \(2N\). Since \(J\) is ample,
\(J^2=2h^0(J)\ge2N\), while Hodge index gives
\(J^2\le(J\cdot H)^2/H^2=2N\). Equality in Hodge index
gives \(J\equiv H\). Numerical and algebraic equivalence
coincide for line bundles on an abelian variety: the numerically
trivial classes give torsion in the Neron--Severi group, and
that group is torsion free by
\citep[Corollaries~7.23 and~7.25, pp.~106--107]
{EdixhovenVanDerGeerMoonenAV}. Thus
\(J\otimes H^{-1}\in\operatorname{Pic}^0(B)\).
Only the ample line \(J\) is needed for this equality argument;
no ampleness of \(D\) in its isotropic case is asserted.

Twisting a line bundle by \(\operatorname{Pic}^0\) changes
neither its translation kernel nor its theta commutator on that
kernel. Indeed the theta extension of a homogeneous line bundle
is commutative, and tensor product multiplies commutators.
Consequently, on \(F(K(L))\),
\[
             e_J(Fx,Fy)=e_H(Fx,Fy)=e_L(x,y)^2.
 \tag{DE-OE9}
\]
But the actual weight-one theta action on \(J\) gives
\(e_J(Fx,Fy)=e_L(x,y)\). Therefore every value of the
nondegenerate theta pairing satisfies \(e_L(x,y)^2=e_L(x,y)\),
and hence equals one. Since \(|K(L)|=N^2>1\), this
contradicts nondegeneracy. The rank-three case is impossible,
and the evaluation rank is four.
\end{proof}

\begin{remark}[Exact scope]
\label{rem:DE-odd-Cartier-evaluation-scope}
\status{Candidate; scope accompanies the unpromoted assertions.}
These statements retain ampleness, global generation and odd
\(N\) as explicit assumptions. The additional generic-generation
assertion requires \(3\nmid N\), and concerns the rank-four
bundle \(F_*L\) on the surface. It does not say that this
bundle is globally generated at every point. Cartier surjectivity
concerns the total map \(\bigwedge^2W\to H^0(H)\); it does
not imply maximal rank of every or of a generic multiplication
map \(\beta(s,-)\), maximal rank of the dual alternating pencil,
smooth-member existence, or full geometric monodromy. In
particular no general type-\((1,\ell)\) full-image theorem or
degree-seventeen conclusion is inferred here.

No finite experiment is needed for these arguments. The full
stability and elliptic-divisor proofs are the previously audited
primitive type-\((1,3)\) dependencies, with the divisor degree
rederived above for arbitrary \(N\). The odd theta-group
representation and central-weight arguments are written explicitly.
\end{remark}
